\documentclass[11pt,a4paper]{article}

\usepackage[spanish,es-noshorthands]{babel}
\usepackage[utf8]{inputenc}
\usepackage[T1]{fontenc}
\usepackage[margin=2.6cm]{geometry}
\usepackage{xcolor}
\usepackage{titlesec}
\usepackage{parskip}
\usepackage{enumitem}
\usepackage{hyperref}

\definecolor{accentblue}{HTML}{1F4E79}
\definecolor{alertred}{HTML}{C0392B}

\hypersetup{colorlinks=true, linkcolor=accentblue, urlcolor=accentblue}

\titleformat{\section}
  {\normalfont\Large\bfseries\color{accentblue}}{}{0pt}{}
\titlespacing*{\section}{0pt}{1.4em}{0.8em}
\newcommand{\seccionrule}{\vspace{-0.6em}{\color{accentblue}\hrule height 0.8pt}\vspace{0.8em}}

\newcommand{\diapo}[2]{%
  \par\medskip\noindent%
  \textbf{\textcolor{accentblue}{#1.}}\enspace #2\par\smallskip}

\setlength{\parindent}{0pt}

\title{\textbf{Guion de la defensa}\\[4pt]
  \large Statistical Properties of the Sharpe Ratio Estimator}
\author{Fabrizio Ortega Suni}
\date{Mayo de 2026}

\begin{document}

\begin{center}
  {\LARGE\bfseries\textcolor{accentblue}{Guion de la defensa}}\\[6pt]
  {\large Statistical Properties of the Sharpe Ratio Estimator}\\[2pt]
  {\normalsize Fabrizio Ortega Suni \quad--\quad Mayo de 2026}
\end{center}

\vspace{0.6em}
\seccionrule

% =====================================================================
\section*{How much can you trust a Sharpe ratio?}
% =====================================================================

\diapo{1}{Buenos días.
Empecemos con un escenario muy simple: supongamos que nos dicen que el
ratio de Sharpe de una estrategia de inversión es de 0.5. La pregunta que
quiero que tengamos en mente los próximos veinte minutos es: ¿qué tan
fiable es este número? Por supuesto, el ratio de Sharpe en sí mismo se
define como el ratio de dos momentos poblacionales, pero en la práctica
se tiene que estimar a partir de una muestra finita, de ahí que haya incertidumbre
alrededor de la estimación.}

\diapo{2}{
El teorema central del límite nos dice que el error de estimación
escalado converge a una normal con cierta varianza. El primer trabajo en
tratar esto rigurosamente es el de Lo (2002), quien obtiene esta varianza
bajo el supuesto de que los rendimientos son IID normales.
Notemos que, una vez tenemos esta varianza asintótica, podemos construir
un intervalo de confianza asintótico alrededor del verdadero SR. Con esto
podemos reportar, con cierto nivel de confianza, dónde se encuentra el
verdadero SR, lo cual es más robusto que solo la estimación puntual.}

\diapo{3}{El problema viene de que la fórmula de Lo depende fuertemente
de la normalidad asumida. Como sabemos de sobra, los retornos financieros
poseen atributos que los alejan de este supuesto, entre ellos colas
pesadas, asimetría, clusters de volatilidad y autocorrelación. Más
adelante veremos cómo se comportan los intervalos y, más generalmente, la
inferencia, cuando tenemos en cuenta o cuando obviamos estos atributos.}

\diapo{4}{Distintos autores han extendido el trabajo de Lo; en esta tabla
muestro los principales. Cada uno de ellos ha asumido un modelo distinto
para los rendimientos y ha obtenido la fórmula de la varianza asintótica
correspondiente. Podemos relacionar cada avance con lo que conté antes:
Mertens incorpora las colas pesadas y la asimetría, y los siguientes, la
autocorrelación y los clusters de volatilidad.
Notemos, sin embargo, que hace falta un resultado para cuando hay
autocorrelación y volatilidad condicional simultáneamente.}

\diapo{5}{Y esta es la principal contribución de este trabajo.
Aquí podemos ver la expresión final; vemos que su presentación ha
requerido múltiples variables auxiliares. Y, por supuesto, los resultados
anteriores se recuperan como casos especiales de esta fórmula.}

\diapo{6}{Volvemos a la diapositiva 4. Notemos que la fórmula IID no
normal permite el siguiente análisis: mientras más grande el SR real, o
más grande la kurtosis, y más negativa la asimetría de los retornos,
mayor será la varianza del estimador (en comparación con la IID normal). Un análisis similar se podría hacer
para el parámetro de autocorrelación phi en el modelo AR(1); sin embargo,
para los demás resulta más difícil, ya que hay interacciones más
complejas entre los parámetros.
Por esto, el trabajo contempla dos análisis de sensibilidad: uno en el
que tomamos un punto base y numéricamente vemos cómo se comporta la
varianza si variamos los demás parámetros, y otro análisis independiente
de un punto base.
Las conclusiones son que los parámetros SR, phi, kurtosis y asimetría
presentan influencias claras, mientras que la influencia de alpha y beta
es conjunta.}

\vspace{0.4em}
\seccionrule

% =====================================================================
\section*{Do the confidence intervals actually work?}
% =====================================================================

\diapo{7--8}{Volviendo a los intervalos, recordemos que, con las
fórmulas anteriores, obtenemos intervalos asintóticos. El trabajo
contiene varios experimentos acerca de qué pasa con estos cuando nos
equivocamos especificando un modelo para los retornos, o cuando,
haciéndolo correctamente, los intervalos no se comportan de buena
manera.
Los experimentos se basan en el concepto de cobertura. Por ejemplo, si
calculamos un intervalo al 95\% de confianza teórico y simulamos 1000
trayectorias Monte Carlo, y contamos cuántas han caído dentro del
intervalo, este número debe ser cercano a 950, o sea, el 95\%.
Ahora veremos unos resultados interesantes. En el primero se generan
sintéticamente procesos AR(1) y se asume, también correctamente, un
modelo AR(1); calculamos la cobertura para distintos valores de
autocorrelación y vemos cómo evolucionan a medida que el tamaño de la
trayectoria crece. Todos convergen al valor esperado, pero vemos que,
cuando la autocorrelación es positiva, para series cortas hay
infracobertura, es decir, nuestro intervalo, dado por la fórmula, se
queda corto: es muy estrecho.
El segundo experimento utiliza series sintéticas de un proceso
AR(1)-GARCH(1,1) y evaluamos nuevamente las coberturas para distintos
tamaños de trayectoria, pero esta vez cada linea representa un distinto modelo asumido. Vemos
que solo uno de los modelos, el AR(1)-GARCH(1,1), converge al valor
esperado; esto quiere decir que cualquier otro modelo —los de expresiones
más simples— infraestima la varianza. Este es un ejemplo de por qué la
fórmula que encontré es necesaria.}

\diapo{9}{Estos son otros dos experimentos. El primero muestra el
escenario de una especificación correcta (GARCH), pero para distintos
niveles de los parámetros alpha+beta. Vemos que, mientras nos acercamos
al punto de no estacionariedad alpha+beta=1, la varianza es sobreestimada
cada vez más, llegando al punto de requerir trayectorias larguísimas para
converger. En la práctica, esto se traduce en que, si sospechamos niveles
altos de persistencia, sabemos que la fórmula de la varianza, aunque
asintóticamente correcta, no es útil.
El segundo nos lleva a otro extremo, esta vez para un proceso IID pero
con innovaciones t de Student; lo que variamos aquí son los grados de
libertad. Vemos que, mientras más pequeños, peor es la convergencia,
llegando al punto de que, para df=3, hay divergencia. Esto es porque,
como se explicó en el trabajo, una de las condiciones de regularidad
necesarias es la existencia del momento de orden cuarto para los
rendimientos, que en este caso no existiría.}

\vspace{0.4em}
\seccionrule

% =====================================================================
\section*{Is the point estimate even right?}
% =====================================================================

\diapo{10}{Dejando de lado los intervalos y volviendo a la estimación
puntual, otra característica del SR sobre la que nos ocupamos es que, en
muestras finitas, es sesgado. Para ello, encontramos las expresiones del
sesgo y su correspondiente corrección para los modelos GARCH y AR-GARCH,
como vemos en la tabla.}

\diapo{11}{Aquí vemos estas correcciones aplicadas a distintos modelos y
vemos cómo aceleran la convergencia (de la estimación puntual). El experimento es muy sencillo, generamos una serie sintética cuyo SR teórico es 0.5 y estimamos puntualmente SR con y sin corrección de sesgo, hacemos esto varias veces y tomamos la media.}

\vspace{0.4em}
\seccionrule

% =====================================================================
\section*{Does the framework survive real data?}
% =====================================================================

\diapo{12}{Hasta ahora solo hemos visto resultados teóricos o su
aplicación a entornos simulados; ahora veremos cómo esto se aplicaría a
un entorno real. Para ello, he tomado tres series semanales de ETF; las
tres son significativas en autocorrelación y heterocedasticidad
condicional, por lo que el modelo AR-GARCH parece un candidato apropiado.
Aquí vemos un gráfico con sus precios y sus rendimientos.}

\diapo{13}{Una vez elegido el modelo, se puede calcular el PSR, que no es
más que uno menos el p-valor del test estadístico cuya hipótesis es
averiguar si el SR es mayor que cierto umbral, en este caso 0.
Intuitivamente, podemos verlo como la probabilidad de que el SR sea mayor
que 0, algo deseable en una inversión.
La tabla muestra el PSR para cada ETF, y las columnas muestran cómo
cambiaría para cada modelo asumido (vemos que la expresión del PSR
depende de la varianza asintótica, cuya fórmula, como vimos al inicio,
cambia dependiendo del modelo asumido).
Lo que quiero mostrar en esta tabla es que la decisión estadística de
aceptar que un SR es positivo cambia dependiendo del modelo que
asumamos. Recordemos que el modelo correcto —por lo menos, para mi
criterio, es el AR-GARCH— pero, digamos que, por ejemplo, un analista
asume un modelo GARCH: estadísticamente, al 5\%, rechazaríamos que
SR$>$0, pero si asumiera el AR-GARCH, sí que aceptaría que SR$>$0.}

\vspace{0.4em}
\seccionrule

% =====================================================================
\section*{Beyond a single number: pairing with VaR}
% =====================================================================

\diapo{14}{Ahora veremos una de las secciones complementarias realizadas en el trabajo. Y es que una crítica posible al uso del ratio de Sharpe es que está
basado en los dos primeros momentos; como alternativa, existen otras
métricas, como el VaR. Imaginemos el escenario de dos inversiones con el mismo SR estimado, pero una de ellas tiene un VaR mayor, obviamente la que tiene menos VaR es más deseable, y viendo solo el SR estaríamos periendo esta información adicional. De aquí surge la idea de que en vez de sustituir este ratio, podemos
trabajar con ambas de forma conjunta; para ello, en este trabajo,
encontramos la distribución asintótica conjunta. Y, dentro de este marco,
analizamos dos alternativas para el estimador del VaR: una completamente
no paramétrica (cuantil empírico) y otra paramétrica (que usa la
propiedad de traslación de los cuantiles).
Lo que encontramos en la matriz de covarianzas asintótica de ambas
opciones es que la covarianza entre ambos estimadores (el SR y el VaR) es
común y negativa, lo que quiere decir que una sobreestimación del SR
suele venir con una infraestimación del VaR. También vemos que el único
término no común entre estas matrices es el de la varianza del VaR. La
alternativa paramétrica es más precisa, y esta diferencia de precisión crece mientras
más pequeño sea el percentil usado para el VaR. Esto nos lleva a pensar que usar el
segundo estimador es más eficiente y que debería ser la elección por defecto, a continuación veremos que no es así.}

\diapo{15}{Para poner a prueba las dos opciones, realizamos nuevamente
experimentos de cobertura, solo que esta vez son dos variables, por lo
que nuestros intervalos serán elipses.
Tenemos dos escenarios: en el primero, los datos (generados) vienen de una normal, y
en el segundo, de una t de Student. En el primero, nuestra identificación
del modelo sería correcta, por lo que, efectivamente, vemos que la
segunda alternativa es más eficiente, logrando una cobertura perfecta del
95\%. En el segundo escenario, habríamos cometido un error de
especificación, y vemos que la alternativa uno, el estimador no
paramétrico, no cambia su cobertura, mientras que el paramétrico se
vuelve sesgado y por lo tanto ya no es fiable.}

\diapo{16}{Entonces, volvemos a nuestra pregunta inicial: ¿podemos
confiar en el SR? La respuesta dependerá de los siguientes
factores:
\begin{enumerate}[label=\arabic*),leftmargin=1.6em,itemsep=2pt,topsep=4pt]
  \item El modelo asumido: la estimación puntual no cambia, pero los
    intervalos sí; y si nos equivocamos en el modelo, estos no serán
    fiables.
  \item Incluso si acertamos con el modelo, la naturaleza de los
    rendimientos —como el signo de su autocorrelación, o si su cuarto
    momento teórico existe— podría hacer que los intervalos tampoco
    sean fiables.
  \item En muestras pequeñas es sesgado, y debemos asegurarnos de que la
    estimación puntual esté corregida.
\end{enumerate}}

\vspace{1.2em}
\seccionrule


\end{document}